\begin{graphillustration}{Слева: частичный порядок (стрелка от меньшего к большему) и его граф сравнимости. Цвет $c(u)$ — число элементов наибольшей цепи, кончающейся в $u$ (цифра у вершины); цепь $a<c<e$ — клика размера $\omega=3$. Справа: «сеть» — не граф сравнимости.}{fig:13-comparability}
\begin{tikzpicture}[font=\small]
\begin{scope}
\node[gi vertex,label=left:{\scriptsize 1}] (a) at (0,0) {$a$};
\node[gi vertex,label=right:{\scriptsize 1}] (b) at (1.6,0) {$b$};
\node[gi vertex,label=left:{\scriptsize 2}] (c) at (0,1.3) {$c$};
\node[gi vertex,label=right:{\scriptsize 2}] (d) at (1.6,1.3) {$d$};
\node[gi vertex,label=left:{\scriptsize 3}] (e) at (0.8,2.6) {$e$};
\draw[gi arrow] (a) -- (c); \draw[gi arrow] (a) -- (d); \draw[gi arrow] (b) -- (d);
\draw[gi arrow] (c) -- (e); \draw[gi arrow] (d) -- (e);
\draw[gi arrow] (a) to[bend left=35] (e);
\draw[gi arrow] (b) to[bend right=35] (e);
\node[gi note] at (0.8,-1.2) {транзитивная ориентация};
\end{scope}
\begin{scope}[shift={(6,0.3)}]
\node[gi vertex] (a) at (0,0) {$a$};
\node[gi vertex] (b) at (1.8,0) {$c$};
\node[gi vertex] (c) at (0.9,1.5) {$b$};
\node[gi vertex] (a1) at (-1,-0.7) {$a'$};
\node[gi vertex] (b1) at (2.8,-0.7) {$c'$};
\node[gi vertex] (c1) at (0.9,2.7) {$b'$};
\draw[gi arrow] (a) -- (c); \draw[gi arrow] (c) -- (b); \draw[gi arrow] (a) -- (b);
\draw[gi edge] (a) -- (a1); \draw[gi edge] (b) -- (b1); \draw[gi edge] (c) -- (c1);
\node[gi note] at (0.9,-1.5) {ребро $bb'$ не ориентировать};
\end{scope}
\end{tikzpicture}
\end{graphillustration}
